\subsection{是否滑动？}

对于存在静摩擦力的平衡系统，求解平衡问题时有一个需要注意的地方就是静摩擦力
并不是可以无限增加的，而是有一最大值。通常来说，静摩擦力的最大值与相互接触且有相对
滑动趋势的物体之间接触面的性质以及相互之间的压力有关，可写成
\[
  f_{\max}=\mu_{\rm s}N,
\]
其中 $\mu_{\rm s}$ 为静摩擦系数。

在处理此类问题时，一般来说首先假设两者间存在静摩擦力，在静摩擦力和其它所有力
共同作用下保持平衡。但最后还要判断平衡所需的摩擦力是否超过上限：若
$|f|\le\mu_{\rm s}N$，则可能保持静止，其中等号对应临界静止；若
$|f|>\mu_{\rm s}N$，假设的静力平衡才不可能实现，物体将发生相对滑动。

\sourcefigure[0.43\linewidth]{fig-1-2.png}
\origcaption{图 1.2：全反力示意图}
图中将接触力矢量平移到共同起点以说明合成关系，不表示这些力实际作用在物块顶面。

除了用摩擦力和压力比值的方法判断是否滑动以外，还可以考虑可能滑动物体与接触
表面之间压力和静摩擦力的合力，有时将其称为全反力。如图 1.2 所示的 $F$ 就是全反力，它
与接触面垂线的夹角 $\alpha$ 是标志全反力的一个重要的物理量。从图中可以看出，当 $\alpha$ 为已知
时，压力 $N=F\cos\alpha$，摩擦力 $f=F\sin\alpha$，所以不发生滑动的条件此时就变成了
\begin{equation}
  \sin\alpha\le\mu_{\rm s}\cos\alpha,\qquad \text{或}\qquad \tan\alpha\le\mu_{\rm s},
\end{equation}
也就是说，全反力和法线夹角的正切不大于静摩擦系数时，接触面可能保持相对静止；等号
对应摩擦角达到极限，超过这一角度才不可能相对静止。

\begin{origexamplebox}{例 1.18}
木箱质量为 $M$，与地面间的动摩擦因数为 $\mu$，用斜向上的力 $F$ 拉木箱使之沿水平
地面匀速前进，如图所示，问角 $\alpha$ 为何值时拉力 $F$ 最小？这个最小值为多大？
\sourcefigure[0.42\linewidth]{ex-1-18.png}
\end{origexamplebox}

\begin{xesolution}
木箱受拉力 $F$、重力 $Mg$、地面支持力 $N$ 和滑动摩擦力 $f$，方向如解图所示：
\soldiagOneEighteen
由解图可见，木箱水平向右匀速运动，摩擦力向左，且 $f=\mu N$。水平方向平衡给出
\[
  F\cos\alpha=\mu N.
\]
竖直方向上，拉力的向上分量减小了地面支持力，故
\[
  N+F\sin\alpha=Mg.
\]
消去 $N$，得到
\[
  F=\frac{\mu Mg}{\cos\alpha+\mu\sin\alpha}.
\]
要使 $F$ 最小，应使分母最大。写成
\[
  \cos\alpha+\mu\sin\alpha
  =\sqrt{1+\mu^2}\cos(\alpha-\arctan\mu),
\]
可知
\[
  \boxed{\alpha=\arctan\mu},
  \qquad
  \boxed{F_{\min}=\frac{\mu Mg}{\sqrt{1+\mu^2}}}.
\]
\end{xesolution}

\begin{origexamplebox}{例 1.19}
筷子夹鸡蛋（只考虑在水平光滑桌面上）的摩擦紧锁问题：设筷子与鸡蛋的摩擦
系数为 $\mu$，求筷子多大角度时在水平桌面上始终不会把鸡蛋滑出。
\sourcefigure[0.44\linewidth]{ex-1-19.png}
\end{origexamplebox}

\begin{xesolution}
取关于角平分线对称的两个接触点。设每根筷子对鸡蛋的正压力为 $N$。两正压力沿角平分
线向外，而极限摩擦力沿筷子指向 $O$：
\soldiagOneNineteen
由解图可知，两正压力沿角平分线向外的合力为
\[
  2N\sin\frac\theta2.
\]
鸡蛋将要向外滑出时，两处摩擦力均沿筷子指向 $O$，其沿角平分线向内的最大合力为
\[
  2\mu N\cos\frac\theta2.
\]
发生摩擦紧锁要求后者不小于前者，即
\[
  2\mu N\cos\frac\theta2\ge 2N\sin\frac\theta2.
\]
所以
\[
  \boxed{\tan\frac\theta2\le\mu},
  \qquad\text{即}\qquad
  \boxed{\theta\le2\arctan\mu}.
\]
\end{xesolution}

\begin{origexamplebox}{例 1.20}
有一长 $l$ 质量为 $M$ 的均匀杆 $AB$，$A$ 端顶在竖直的粗糙墙壁上，杆端与墙间的摩擦
系数为 $\mu$，$B$ 端用一强度足够而不可伸长的绳悬挂，绳的另一端固定在墙壁 $C$ 点，木杆呈
水平状态，绳与杆的夹角为 $\theta$，求杆能保持平衡时 $\mu$ 与 $\theta$ 应满足的条件。杆保持平衡时，杆
上有一点 $P$ 存在，若 $A$ 与 $P$ 点间挂一质量足够大的重物，则可以破坏平衡，而在 $PB$ 间
任一点悬挂任意重物均不能破坏平衡。求 $PA$ 距离 $x$。
\sourcefigure[0.46\linewidth]{ex-1-20.png}
\end{origexamplebox}

\begin{xesolution}
杆、墙面和细绳的受力关系如解图所示；附加重物的位置暂记为 $y$：
\soldiagOneTwenty
先考虑没有附加重物的木杆。设绳张力为 $T$，墙面对杆的支持力和摩擦力分别为 $N$、
$f$。对 $A$ 点取矩：
\[
  T\sin\theta\,l=Mg\frac l2,
  \qquad
  T\sin\theta=\frac{Mg}{2}.
\]
再由水平、竖直方向平衡得到
\[
  N=T\cos\theta=\frac{Mg}{2}\cot\theta,
  \qquad
  f=Mg-T\sin\theta=\frac{Mg}{2}.
\]
静摩擦条件 $f\le\mu N$ 因而给出
\[
  \boxed{\mu\ge\tan\theta}.
\]

现在在距 $A$ 点为 $y$ 的位置悬挂质量为 $m$ 的重物。重复上述计算：
\[
  T\sin\theta=\frac{Mg}{2}+mg\frac yl,
\]
\[
  f=\frac{Mg}{2}+mg\left(1-\frac yl\right),
  \qquad
  N=\left(\frac{Mg}{2}+mg\frac yl\right)\cot\theta.
\]
若无论 $m$ 多大都不能破坏平衡，则除原先的 $\mu\ge\tan\theta$ 外，还必须使 $m$ 的
系数满足
\[
  1-\frac yl\le\mu\frac yl\cot\theta.
\]
临界点 $P$ 满足等号，于是
\[
  \boxed{x=AP=\frac{l\tan\theta}{\mu+\tan\theta}}.
\]
当悬挂点位于 $AP$ 内时，可用足够重的物体破坏平衡；位于 $PB$ 内时，任意重物都不能
使接触处超过静摩擦极限。
\end{xesolution}

\begin{origexamplebox}{例 1.21}
如图，一人对一均匀细杆的一端施力，力的方向总与杆垂直，要将杆从地板上无
滑动地慢慢抬到竖直位置，试求杆与地板间的静摩擦因数至少应为多大？
\sourcefigure[0.40\linewidth]{ex-1-21.png}
\end{origexamplebox}

\begin{xesolution}
设杆长为 $L$、质量为 $m$，杆与地板夹角为 $\theta$。外力 $F$ 始终垂直于杆。对杆的下端
取矩最方便，因为地板处的两个力都不产生力矩：
\soldiagOneTwentyOne
由解图有
\[
  FL=mg\frac L2\cos\theta,
  \qquad
  F=\frac{mg}{2}\cos\theta.
\]
由水平、竖直方向平衡可得地板处的摩擦力和支持力：
\[
  f=F\sin\theta=\frac{mg}{2}\sin\theta\cos\theta,
\]
\[
  N=mg-F\cos\theta
  =\frac{mg}{2}\left(1+\sin^2\theta\right).
\]
整个抬起过程中均不滑动，要求
\[
  \mu\ge\max_{0\le\theta\le\pi/2}\frac fN
  =\max_{0\le\theta\le\pi/2}
  \frac{\sin\theta\cos\theta}{1+\sin^2\theta}.
\]
在 $0\le\theta<\pi/2$ 内令 $t=\tan\theta\ge0$，上式化为 $t/(1+2t^2)$。
由非负平方 $(\sqrt2t-1)^2\ge0$ 得
\[
 1+2t^2\ge2\sqrt2t.
\]
当 $t>0$ 时，两边除以正数 $1+2t^2$，得
$t/(1+2t^2)\le1/(2\sqrt2)$，等号恰在 $t=1/\sqrt2$ 时成立。
$t=0$ 时此比值为零；当 $\theta\to\pi/2$、$t\to\infty$ 时也趋于零，
所以两端都不提供更大的摩擦要求。因此
\[
  \boxed{\mu_{\min}=\frac{1}{2\sqrt2}}.
\]
\end{xesolution}

\begin{origexamplebox}{例 1.22}
质量均为 $m$ 的两环 $A$、$B$ 用长为 $a$ 的细线相连在水平杆上，在细线的中点拴有
一质量为 $M$ 的物块 $C$，$A$、$B$ 环与杆间的静摩擦系数为 $\mu$，求平衡情况下的两环的最大距离
$x$。
\sourcefigure[0.50\linewidth]{ex-1-22.png}
\end{origexamplebox}

\begin{xesolution}
系统关于中点对称。每一侧细线长度为 $a/2$。设细线与水平方向夹角为 $\beta$，则
\[
  \cos\beta=\frac xa,
  \qquad
  \sin\beta=\frac{\sqrt{a^2-x^2}}{a}.
\]
中间结点和任一圆环的受力应分开研究：
\soldiagOneTwentyTwo
在中间结点处列竖直方向平衡方程：
\[
  2T\sin\beta=Mg,
  \qquad
  T=\frac{Mg}{2\sin\beta}.
\]
对任一圆环，细线的竖直分量为 $T\sin\beta=Mg/2$，所以杆对环的支持力为
\[
  N=\left(m+\frac M2\right)g.
\]
环将要滑动时
\[
  T\cos\beta=\mu N.
\]
代入 $T$ 和几何关系得
\[
  Mx=\mu(2m+M)\sqrt{a^2-x^2}.
\]
因此两环能够保持平衡的最大距离为
\[
  \boxed{x_{\max}=
  \frac{a\mu(2m+M)}{\sqrt{M^2+\mu^2(2m+M)^2}}}.
\]
\end{xesolution}
